% Tablas 2--12 del capítulo IV, construidas a partir del Excel de tablas aprobado.
% El título y el número son automáticos mediante captionof{table}.
% Tabla 1 corresponde a la matriz del capítulo I.

\newcommand{\tablaResultadoDos}{%
\par\medskip\noindent
\begingroup
\setlength{\parskip}{0pt}
\setlength{\tabcolsep}{4pt}
\footnotesize
\captionof{table}{Distribución de unidades estructura–sector por tramo y condición superficial}\label{tab:res-distribucion}
\vspace{4pt}
\resizebox{\linewidth}{!}{%
\begin{tabular}{@{}lrrrrr@{}}
\toprule
\textbf{Tramo} & \textbf{Unidades} & \textbf{Recapeo} & \textbf{Sector mixto} & \textbf{Pavimento flexible} & \textbf{Porcentaje} \\
\midrule
Tramo 1 & 14 & 2 & 1 & 11 & 70\% \\
Tramo 2 & 2 & 0 & 0 & 2 & 10\% \\
Tramo 3 & 2 & 0 & 0 & 2 & 10\% \\
Tramo 4 & 2 & 0 & 0 & 2 & 10\% \\
TOTAL & 20 & 2 & 1 & 17 & 100\% \\
\bottomrule
\end{tabular}
}% fin resizebox
\notaTabla{Datos calculados a partir de la base unificada de las veinte unidades estructura--sector de Simbal, 2026.}
\endgroup
}

\newcommand{\tablaResultadoTres}{%
\begin{paginaHorizontal}
\begingroup
\setlength{\parskip}{0pt}
\setlength{\tabcolsep}{4pt}
\scriptsize
\captionof{table}{Indicadores CL, CT, PCI e IRI de las veinte unidades estructura–sector}\label{tab:res-indicadores}
\vspace{4pt}
\begin{tabularx}{\linewidth}{@{}llrlrrrr@{}}
\toprule
\textbf{Unidad} & \textbf{Tramo} & \textbf{Pase (m)} & \textbf{Superficie} & \textbf{CL (\%)} & \textbf{CT (\%)} & \textbf{PCI} & \textbf{IRI (m/km)} \\
\midrule
SIM-01 & 1 & 42 & Recapeo & 74,20 & 78,64 & 78,50 & 4,6175 \\
SIM-02 & 1 & 162 & Recapeo & 83,47 & 78,35 & 83,39 & 2,3875 \\
SIM-03 & 1 & 295 & Sector mixto & 67,92 & 79,85 & 60,93 & 3,8025 \\
SIM-04 & 1 & 443 & Pav. flexible & 89,19 & 66,44 & 80,45 & 2,2225 \\
SIM-05 & 1 & 611 & Pav. flexible & 83,50 & 84,40 & 82,01 & 1,9200 \\
SIM-06 & 1 & 802 & Pav. flexible & 74,72 & 70,39 & 73,85 & 2,5625 \\
SIM-07 & 1 & 1009 & Pav. flexible & 92,91 & 78,43 & 79,86 & 3,0875 \\
SIM-08 & 1 & 1183 & Pav. flexible & 90,92 & 87,10 & 64,48 & 2,1575 \\
SIM-09 & 1 & 1385 & Pav. flexible & 70,70 & 72,43 & 72,49 & 4,0675 \\
SIM-10 & 1 & 1584 & Pav. flexible & 74,94 & 65,74 & 60,06 & 2,6575 \\
SIM-11 & 1 & 1795 & Pav. flexible & 63,89 & 77,20 & 71,95 & 2,6775 \\
SIM-12 & 1 & 1982 & Pav. flexible & 70,03 & 84,65 & 62,99 & 2,7850 \\
SIM-13 & 1 & 2225 & Pav. flexible & 63,77 & 91,26 & 64,31 & 4,1450 \\
SIM-14 & 1 & 2477 & Pav. flexible & 55,88 & 71,94 & 50,26 & 4,2525 \\
SIM-15 & 2 & 33 & Pav. flexible & 42,50 & 74,66 & 67,45 & 3,2600 \\
SIM-16 & 2 & 108 & Pav. flexible & 49,56 & 69,54 & 65,90 & 3,0825 \\
SIM-17 & 3 & 65 & Pav. flexible & 38,50 & 61,71 & 51,35 & 4,7175 \\
SIM-18 & 3 & 235 & Pav. flexible & 88,59 & 83,30 & 75,86 & 3,1200 \\
SIM-19 & 4 & 34 & Pav. flexible & 76,02 & 85,79 & 64,96 & 3,1425 \\
SIM-20 & 4 & 182 & Pav. flexible & 70,25 & 67,19 & 62,41 & 3,9025 \\
\bottomrule
\end{tabularx}
\notaTabla{Datos calculados a partir de la base unificada de las veinte unidades estructura--sector de Simbal, 2026.}
\endgroup
\end{paginaHorizontal}
}

\newcommand{\tablaResultadoCuatro}{%
\par\medskip\noindent
\begingroup
\setlength{\parskip}{0pt}
\setlength{\tabcolsep}{4pt}
\footnotesize
\captionof{table}{Verificación de reconstrucción de indicadores y observaciones de consistencia de la base}\label{tab:res-auditoria}
\vspace{4pt}
\begin{tabularx}{\linewidth}{@{}p{2.7cm}p{1.35cm}p{2.1cm}X p{1.7cm}@{}}
\toprule
\textbf{Indicador / unidad} & \textbf{Registros} & \textbf{Diferencia} & \textbf{Comprobación} & \textbf{Resultado} \\
\midrule
CL & 20 & 0 & \parbox[t]{5.5cm}{Cálculo CL reconstruido desde Lc y Le} & Conforme \\
CT & 20 & 0 & \parbox[t]{5.5cm}{Cálculo CT reconstruido desde áreas de ambos extremos} & Conforme \\
PCI & 20 & 0 & \parbox[t]{5.5cm}{Promedio sectorial ponderado por área} & Conforme \\
IRI & 20 & 0 & \parbox[t]{5.5cm}{Consolidación de los recorridos de ambas huellas} & Conforme \\
SIM-14 & 1 & 0,21 & \parbox[t]{5.5cm}{Longitud nominal 82.46 m; perfilada 82.25 m} & Observación \\
SIM-16 & 1 & 0,11 & \parbox[t]{5.5cm}{Longitud nominal 84.11 m; perfilada 84.00 m} & Observación \\
SIM-20 & 1 & 0,24 & \parbox[t]{5.5cm}{Longitud nominal 85.74 m; perfilada 85.50 m} & Observación \\
\bottomrule
\end{tabularx}
\notaTabla{Los errores numéricos de reconstrucción se redondean a cero a la precisión mostrada; las diferencias longitudinales se expresan en metros.}
\endgroup
}

\newcommand{\tablaResultadoCinco}{%
\begin{paginaHorizontal}
\begingroup
\setlength{\parskip}{0pt}
\setlength{\tabcolsep}{4pt}
\scriptsize
\captionof{table}{Estadísticos descriptivos de CL, CT, PCI e IRI en las veinte unidades}\label{tab:res-descriptivos}
\vspace{4pt}
\begin{tabularx}{\linewidth}{@{}llrrrrrrrrr@{}}
\toprule
\textbf{Indicador} & \textbf{Unidad} & \textbf{n} & \textbf{Media} & \textbf{DE} & \textbf{Mediana} & \textbf{Q1} & \textbf{Q3} & \textbf{RIC} & \textbf{Mín.} & \textbf{Máx.} \\
\midrule
CL & \% & 20 & 71,07 & 15,50 & 72,45 & 63,86 & 83,47 & 19,61 & 38,50 & 92,91 \\
CT & \% & 20 & 76,45 & 8,16 & 77,78 & 70,18 & 83,57 & 13,39 & 61,71 & 91,26 \\
PCI & 0–100 & 20 & 68,67 & 9,63 & 66,68 & 62,84 & 76,52 & 13,68 & 50,26 & 83,39 \\
IRI & m/km & 20 & 3,23 & 0,84 & 3,10 & 2,63 & 3,94 & 1,31 & 1,92 & 4,72 \\
\bottomrule
\end{tabularx}
\notaTabla{Datos calculados a partir de la base unificada de las veinte unidades estructura--sector de Simbal, 2026.}
\endgroup
\end{paginaHorizontal}
}

\newcommand{\tablaResultadoSeis}{%
\par\medskip\noindent
\begingroup
\setlength{\parskip}{0pt}
\setlength{\tabcolsep}{4pt}
\footnotesize
\captionof{table}{Promedios de CL, CT, PCI e IRI según tramo vial}\label{tab:res-promedios}
\vspace{4pt}
\resizebox{\linewidth}{!}{%
\begin{tabular}{@{}lrrrrr@{}}
\toprule
\textbf{Tramo} & \textbf{n} & \textbf{CL media (\%)} & \textbf{CT media (\%)} & \textbf{PCI medio} & \textbf{IRI medio (m/km)} \\
\midrule
1 & 14 & 75,43 & 77,63 & 70,39 & 3,10 \\
2 & 2 & 46,03 & 72,10 & 66,68 & 3,17 \\
3 & 2 & 63,55 & 72,50 & 63,60 & 3,92 \\
4 & 2 & 73,14 & 76,49 & 63,69 & 3,52 \\
\bottomrule
\end{tabular}
}% fin resizebox
\notaTabla{Datos calculados a partir de la base unificada de las veinte unidades estructura--sector de Simbal, 2026.}
\endgroup
}

\newcommand{\tablaResultadoSiete}{%
\par\medskip\noindent
\begingroup
\setlength{\parskip}{0pt}
\setlength{\tabcolsep}{4pt}
\footnotesize
\captionof{table}{Correlaciones de Spearman entre indicadores de drenaje y desempeño funcional}\label{tab:res-spearman}
\vspace{4pt}
\resizebox{\linewidth}{!}{%
\begin{tabular}{@{}llrrrl@{}}
\toprule
\textbf{Hipótesis} & \textbf{Relación} & \textbf{n} & \textbf{Rho} & \textbf{p bilateral} & \textbf{Sentido} \\
\midrule
H1 & CL–PCI & 20 & 0,5940 & 0,0058 & Directa \\
H2 & CL–IRI & 20 & -0,6105 & 0,0042 & Inversa \\
H3 & CT–PCI & 20 & 0,1714 & 0,4699 & Directa \\
H4 & CT–IRI & 20 & -0,1053 & 0,6587 & Inversa \\
\bottomrule
\end{tabular}
}% fin resizebox
\notaTabla{Datos calculados a partir de la base unificada de las veinte unidades estructura--sector de Simbal, 2026.}
\endgroup
}

\newcommand{\tablaResultadoOcho}{%
\par\medskip\noindent
\begingroup
\setlength{\parskip}{0pt}
\setlength{\tabcolsep}{4pt}
\footnotesize
\captionof{table}{Parámetros de ajuste de tendencia espacial e índice de Moran de los rangos y residuos}\label{tab:res-moran-param}
\vspace{4pt}
\resizebox{\linewidth}{!}{%
\begin{tabular}{@{}lrrrrr@{}}
\toprule
\textbf{Indicador} & \textbf{Intercepto} & \textbf{$\beta_X$} & \textbf{$\beta_Y$} & \textbf{Moran rangos} & \textbf{Moran residuos} \\
\midrule
CL & 10,5000 & -2,6107 & 1,3866 & -0,1232 & -0,1053 \\
CT & 10,5000 & 0,5320 & -0,0227 & 0,0266 & 0,0179 \\
PCI & 10,5000 & -4,9297 & 1,5757 & 0,0480 & -0,3506 \\
IRI & 10,5000 & 1,3866 & 0,4364 & 0,0642 & -0,0691 \\
\bottomrule
\end{tabular}
}% fin resizebox
\notaTabla{Datos calculados a partir de la base unificada de las veinte unidades estructura--sector de Simbal, 2026.}
\endgroup
}

\newcommand{\tablaResultadoNueve}{%
\par\medskip\noindent
\begingroup
\setlength{\parskip}{0pt}
\setlength{\tabcolsep}{4pt}
\footnotesize
\captionof{table}{Correlaciones ajustadas y valores p de las 9 999 aleatorizaciones espaciales}\label{tab:res-moran-pruebas}
\vspace{4pt}
\begin{tabularx}{\linewidth}{@{}llrrrrl@{}}
\toprule
\textbf{Hipótesis} & \textbf{Relación} & \textbf{Rho residual} & \textbf{B} & \textbf{b extremos} & \textbf{p espacial} & \textbf{Prueba} \\
\midrule
H1 & CL–PCI & 0,5773 & 9999 & 44 & 0,0045 & Bilateral \\
H2 & CL–IRI & -0,5893 & 9999 & 355 & 0,0356 & Bilateral \\
H3 & CT–PCI & 0,2893 & 9999 & 2674 & 0,2675 & Bilateral \\
H4 & CT–IRI & -0,1400 & 9999 & 5216 & 0,5217 & Bilateral \\
\bottomrule
\end{tabularx}
\notaTabla{El valor p bilateral se calculó como (b + 1)/(B + 1), con B = 9\,999 aleatorizaciones.}
\endgroup
}

\newcommand{\tablaResultadoDiez}{%
\begin{paginaHorizontal}
\begingroup
\setlength{\parskip}{0pt}
\setlength{\tabcolsep}{4pt}
\footnotesize
\captionof{table}{Ajuste secuencial de Holm y decisiones estadísticas de las hipótesis H1 a H4}\label{tab:res-holm}
\vspace{4pt}
\begin{tabularx}{\linewidth}{@{}rlrrrrrl@{}}
\toprule
\textbf{Orden} & \textbf{Hipótesis} & \textbf{p espacial} & \textbf{Factor} & \textbf{p por factor} & \textbf{p Holm} & \textbf{Umbral} & \textbf{Decisión} \\
\midrule
1 & H1 & 0,0045 & 4 & 0,0180 & 0,0180 & 0,0125 & Significativa \\
2 & H2 & 0,0356 & 3 & 0,1068 & 0,1068 & 0,0167 & No significativa \\
3 & H3 & 0,2675 & 2 & 0,5350 & 0,5350 & 0,0250 & No significativa \\
4 & H4 & 0,5217 & 1 & 0,5217 & 0,5350 & 0,0500 & No significativa \\
\bottomrule
\end{tabularx}
\notaTabla{Significancia familiar: $\alpha=0{,}05$; ajuste secuencial de Holm.}
\endgroup
\end{paginaHorizontal}
}

\newcommand{\tablaResultadoOnce}{%
\par\medskip\noindent
\begingroup
\setlength{\parskip}{0pt}
\setlength{\tabcolsep}{4pt}
\footnotesize
\captionof{table}{Intervalos de sensibilidad de las correlaciones de Spearman por exclusión de una unidad}\label{tab:res-loo}
\vspace{4pt}
\resizebox{\linewidth}{!}{%
\begin{tabular}{@{}llrrrlr@{}}
\toprule
\textbf{Hipótesis} & \textbf{Relación} & \textbf{Rho inicial} & \textbf{Mín. LOO} & \textbf{Máx. LOO} & \textbf{Signo estable} & \textbf{Exclusiones} \\
\midrule
H1 & CL–PCI & 0,5940 & 0,5316 & 0,6807 & Sí & 20 \\
H2 & CL–IRI & -0,6105 & -0,6789 & -0,5456 & Sí & 20 \\
H3 & CT–PCI & 0,1714 & 0,0439 & 0,2684 & Sí & 20 \\
H4 & CT–IRI & -0,1053 & -0,2211 & 0,0439 & No & 20 \\
\bottomrule
\end{tabular}
}% fin resizebox
\notaTabla{LOO: exclusión sucesiva de una unidad; las correlaciones se calcularon con $n=19$.}
\endgroup
}

\newcommand{\tablaResultadoDoce}{%
\begin{paginaHorizontal}
\begingroup
\setlength{\parskip}{0pt}
\setlength{\tabcolsep}{4pt}
\footnotesize
\captionof{table}{Síntesis de resultados de los objetivos específicos, hipótesis y decisión inferencial}\label{tab:res-sintesis}
\vspace{4pt}
\begin{tabularx}{\linewidth}{@{}llllrrrl@{}}
\toprule
\textbf{Objetivo} & \textbf{Hipótesis} & \textbf{Relación} & \textbf{Sentido} & \textbf{Rho residual} & \textbf{p espacial} & \textbf{p Holm} & \textbf{Resultado} \\
\midrule
OE1 & H1 & CL–PCI & Directa & 0,5773 & 0,0045 & 0,0180 & Significativa (Holm) \\
OE2 & H2 & CL–IRI & Inversa & -0,5893 & 0,0356 & 0,1068 & No significativa (Holm) \\
OE3 & H3 & CT–PCI & Directa & 0,2893 & 0,2675 & 0,5350 & No significativa (Holm) \\
OE4 & H4 & CT–IRI & Inversa & -0,1400 & 0,5217 & 0,5350 & No significativa (Holm) \\
\bottomrule
\end{tabularx}
\notaTabla{Las decisiones corresponden a hipótesis específicas; no se estimó un contraste global independiente.}
\endgroup
\end{paginaHorizontal}
}
